% Section 1.4, from lectures/L01/appendix/d_first_subroutine.md.

\section{Your First Subroutine}
\label{c1:sec:first}

\subsection{The task}
\label{c1:sec:task}
Write two subroutines in \file{drivers/source/utils.asm}, and one program in
\file{drivers/app/main.asm}.

The subroutines compute \code{1 << n} and \code{~(1 << n)} for a value of \code{n} that is not
known until the program runs. In C that is one operator each and you would not give it a second
thought. Here it is a loop, because the AVR has no barrel shifter (\secref{c1:sec:absent}):
\code{lsl} shifts by exactly one place, and shifting by \code{n} places means going round \code{n}
times.

That is not a curiosity. Every driver in this book is handed a pin number and has to turn it into
a bit mask, and these two subroutines are how. \code{led_on} in Chapter~2 calls
\code{shift_bits}; \code{btn_init} in Chapter~3 calls both. Getting them right now means the rest
of the book is about peripherals rather than about bit twiddling.

Creating the file is all it takes to switch its tests on; the README beside the suite, in
\file{lectures/L01/exercises/test}, says how.

\subsection{Where the arguments live}
\label{c1:sec:args}
This book follows the AVR C calling convention from the first line of assembly it writes, and
keeps following it, so that Chapter~6's ``call your assembly from C'' is a demonstration rather
than a rewrite. The whole of what you need for now:

\begin{booktable}{@{}p{9.4em}L@{}}
  \thead{Register} & \thead{Role} \\
  \midrule
  \code{r24} & First argument, if it is a byte. Also where a byte-sized return value goes. \\
  \code{r25:r24} & First argument, if it is a pointer or another 16-bit value. \\
  \code{r22} & Second argument, if it is a byte. \\
  \code{r18} to \code{r27}, \code{r30}, \code{r31} & \textbf{Call-clobbered.} A subroutine may
    destroy these freely. \\
  \code{r2} to \code{r17}, \code{r28}, \code{r29} & \textbf{Call-saved.} A subroutine that uses
    one must put it back. \\
\end{booktable}

So both subroutines here take their argument in \code{r24} and leave their result in \code{r24},
and both may use \code{r18} and \code{r19} as scratch without saving anything.

\subsection{\code{shift_bits}}
\label{c1:sec:shiftbits}
Returns \code{1} shifted left by the number of places given.

\begin{narrowtable}{@{}ccl@{}}
  \thead{Given \code{r24}} & \thead{Returns \code{r24}} & \thead{Because} \\
  \midrule
  0 & \code{0x01} & Shifted no times at all. \\
  1 & \code{0x02} & \\
  5 & \code{0x20} & Which is bit 5, the Arduino Uno's pin 13. \\
  7 & \code{0x80} & \\
  8 & \code{0x00} & The bit has been shifted off the end of an eight-bit register. \\
\end{narrowtable}

Write it like this:
\begin{enumerate}
  \item Put \code{0x01} in a scratch register; this is the value being shifted. Put \code{0} in a
    second scratch register; this is the count of shifts done so far.
  \item At the top of the loop, compare the count against \code{r24}. If they are equal, leave the
    loop.
  \item Otherwise shift the value left once, add one to the count, and jump back to the top.
  \item On leaving, move the value into \code{r24} and return.
\end{enumerate}

Two things about that shape are worth being deliberate about, because they are what the cost
depends on (\secref{c1:sec:branchcost}):

\textbf{The comparison is at the top, not the bottom.} So an argument of zero does no shifts at
all, which is what makes \code{shift_bits(0)} return \code{0x01} rather than \code{0x02}. A loop
with the test at the bottom always runs its body once, and would be wrong for exactly one input,
which is the input your driver will use for pin 8.

\textbf{The exit branch runs once more than the body does.} That last comparison, the one that
finds the count has reached \code{r24} and leaves, is a real instruction with a real cost. It is
the reason the answer is 10 cycles at $n = 0$ and not 7.

The last row of the table is the interesting one. Nothing in the routine says ``stop at eight'';
the zero falls out of the register being eight bits wide, and the bit simply leaves. A reader
thinking in C, where \code{1 << 8} is 256, will predict \code{0x100} and there is nowhere to put
it.

\subsection{\code{shift_bits_inverted}}
\label{c1:sec:inverted}
Returns the complement of the same value: \code{~(1 << n)}.

\begin{narrowtable}{@{}cc@{}}
  \thead{Given \code{r24}} & \thead{Returns \code{r24}} \\
  \midrule
  0 & \code{0xFE} \\
  1 & \code{0xFD} \\
  5 & \code{0xDF} \\
  7 & \code{0x7F} \\
  8 & \code{0xFF} \\
\end{narrowtable}

The obvious implementation is to call \code{shift_bits} and complement the result with
\code{com}. Do not; write it as its own loop. Start the scratch value at \code{0xFE} instead of
\code{0x01}, and each time round shift left \emph{and then increment}, which shifts a zero in at
the bottom and turns it back into a one.

The reason for the detour is that this routine exists to produce a mask for clearing a bit, and
Chapter~2's driver calls it directly. Writing it as its own loop also makes the point that
\secref{c1:sec:branchcost} is really making: two routines differing by a single \code{inc} differ in
cost by that \code{inc} on every trip, and you can say by how much before running either.

\subsection{Why \code{r18} and \code{r19}}
\label{c1:sec:whyr18}
Use \code{r18} and \code{r19} as the scratch registers, not \code{r16} and \code{r17}.

Both pairs are in the upper half, so both can take an \code{ldi} (\secref{c1:sec:regfile}). The
difference is the calling convention in \secref{c1:sec:args}: \code{r18} and \code{r19} are
call-clobbered, and \code{r16} and \code{r17} are call-saved. A subroutine that writes to
\code{r16} without restoring it has broken its contract with whoever called it, and the caller has
no way to know.

Right now nothing calls these subroutines but a test, and a test does not care. By Chapter~6 the
caller may be a C function that had a live value in \code{r16}, and then the bug is a variable
that changes on its own across a function call, which is about as unpleasant a thing to debug as
this book contains. Choosing the right registers now costs nothing and means that never happens.

If you are comparing against AVR assembly you have found elsewhere, note that a good deal of it
uses \code{r16} and \code{r17} here. It works, right up until it does not.

\subsection{The program}
\label{c1:sec:program}
\file{drivers/app/main.asm} is the first thing in this book that is a program rather than a
subroutine. For this chapter it needs four things:

\textbf{A reset vector.} The first word of program memory is where the machine starts, so the
first instruction in your file has to be a jump to your code. Use \code{.org 0} followed by an
\code{rjmp}.

\textbf{A stack pointer.} \code{rcall} pushes a return address and \code{ret} pops one, and both
use \code{SP}. Set it to \code{RAMEND} by writing \code{high(RAMEND)} to \code{SPH} and then
\code{low(RAMEND)} to \code{SPL}, with \code{out} rather than \code{sts}
(\secref{c1:sec:iowindow}). \code{low()} and \code{high()} are AVRASM2's; GNU \code{as} spells the
same two functions \code{lo8()} and \code{hi8()}, which matters once in Chapter~6 and nowhere
else.

The hardware already does this on the ATmega328P, and it does it after \textbf{every} reset,
watchdog resets included. You will find a good deal of AVR code carrying a comment saying
otherwise; on this device that comment is wrong, and repeating it is how the belief survives.

Write the four instructions anyway, for the two reasons that are actually true. Older AVRs reset
\code{SP} to \code{0x0000} rather than to \code{RAMEND}, so code that omits this is not portable
to them, and the failure there is a stack that grows down out of memory on the first
\code{rcall}. And a reset is not the only way to arrive at your code: a bootloader that jumps into
it has performed no reset at all, and \code{SP} holds whatever the bootloader left in it.

Neither of those applies to anything in this book. Knowing \emph{which} reason applies, rather
than carrying a habit you cannot justify, is the point.

\textbf{Something to do, and somewhere to stop.} Call \code{shift_bits} with an argument, then sit
in a loop that jumps to itself. There is no operating system to return to; a program that runs off
the end of its own code executes whatever is in flash after it, which is \code{0xFF} repeated, and
the results are not interesting. An infinite loop is how an embedded program ends.

\textbf{And something you can see.} Light the LED on PB5.

\begin{asm}
    sbi  DDRB, 5                ; PB5 is an output
    sbi  PORTB, 5               ; and it is high
\end{asm}

That is the whole of it, and every part of it is already in this chapter. \code{DDRB} is
\code{0x04} and \code{PORTB} is \code{0x05}, both I/O addresses, so \code{sbi} reaches them and
\code{sts} would not (\secref{c1:sec:iowindow}); \code{sbi} costs two cycles and not one, which
\secref{c1:sec:costs} says and which is worth noticing now rather than in Chapter~5 when you are
counting a timer's handler.

The one thing that is new is what the two registers mean, and it is one sentence:
\textbf{\code{DDRx} decides whether a pin is an output, and \code{PORTx} decides what an output
drives.} Both are per-bit, both default to zero, and zero means ``input, not driven''. So an LED
needs both writes and in that order: setting \code{PORTB} first would briefly enable a pull-up on
an input pin instead, which is a different thing that happens to look similar.

PB5 is Arduino pin 13, which is the one with an LED already fitted to the board, and the reason
this book lights that one rather than any other. On a real Uno there is a resistor beside it; in
the simulator there is nothing to protect, and \secref{c2:sec:resistor} is where that stops being a detail.

\textbf{Two instructions, one pin, hard-coded.} That is deliberate, and it is exactly as far as
this chapter goes. It only works for PB5, it cannot be told which LED to light, and nothing else
in the program can reuse it. Turning it into something a program can call with a pin number is
Chapter~2, and so is the reason that turns out to be more work than it sounds.

\subsection{How to check it}
\label{c1:sec:check}
Three things, in this order.

\textbf{Does it assemble?}

\begin{shell}
make build
\end{shell}

\textbf{Does it do the right thing, and cost the right amount?}

\begin{shell}
make test
\end{shell}

Nine tests were passing before you wrote anything, eighteen once \file{utils.asm} exists, and
twenty once \file{main.asm} does. If \code{Utils.SubroutinesAreDefined} fails while the file is
plainly there, the label is spelled differently from the name the test asks for, or the file did
not assemble (\secref{c1:sec:avra}).

The tests for \file{main.asm} work differently from the rest, and it is worth knowing why before
one of them fails. A subroutine can be called: a test sets the program counter to
\code{shift_bits}, puts a number in \code{r24} and reads \code{r24} back. A program cannot. It has
no arguments, it never returns, and the only way to find out what it does is to \textbf{let it
run} and then look at the machine:

\begin{cppcode}
mcu.run(2000);                      // 2000 cycles of your program
mcu.data(0x24) & (1 << 5)           // DDRB, in the data space
\end{cppcode}

Note the \code{0x24}. The test reads the \textbf{data space} address where your program wrote the
\textbf{I/O} address \code{0x04}, and both are correct, because they are two names for one
register. That is \secref{c1:sec:iowindow} again, seen from the third side, and it is the last time
this book will point it out.

\textbf{What does it actually cost?}

\begin{shell}
make measure SYMBOL=shift_bits ARG=5
\end{shell}

Do not run that until you have predicted the answer. Predicting it by hand, and then reconciling
the prediction against the measurement, is the cross-check in Exercise~\ref{c1:ex:crosscheck}, and
it is the exercise this whole chapter exists to set up.
